% !TEX root = ../../../main.tex
\section{Pullbacks and volume functions}
\label{sec:linear-algebra-i:pullbacks-and-volume-functions}

% Sources: IMG_9262--IMG_9263; October 8 PDF, pp. 13--14.
% The photo uses X* for scalar-valued functions on a set. Use F(X,F) here
% to distinguish this general function space from the linear dual V*.
\subsection{Pullback by composition}

Let \(X,Y\) be sets and let \(F:X\to Y\) be a map. Write
\(\mathcal F(X,\mathbb F)\) for the scalar-valued functions on \(X\).
Every function \(f:Y\to\mathbb F\) gives a function on \(X\) by composition,
\[
  F^*f=f\circ F,\qquad (F^*f)(x)=f(F(x)).
\]
This construction is called pullback. Its direction is opposite to that of
the original map,
\[
  F:X\longrightarrow Y,\qquad
  F^*:\mathcal F(Y,\mathbb F)\longrightarrow\mathcal F(X,\mathbb F).
\]

\autoref{fig:linear-algebra-i:pullback-composition} separates the path
of an input from the direction in which a function is pulled back.
\begin{marginfigure}
  \centering
  % Pullback evaluates by following F and then f; the lower arrow acts on functions.
\begin{tikzpicture}[
  x=1mm,y=1mm,font=\footnotesize,line width=0.45pt,
  every node/.style={inner sep=3pt},
  map/.style={-{Latex[length=1.6mm,width=1.1mm]}}
]
  \node (X) at (0,26) {\(X\)};
  \node (Y) at (30,26) {\(Y\)};
  \node (field) at (30,3) {\(\mathbb F\)};
  \draw[map] (X) -- node[above] {\(F\)} (Y);
  \draw[map] (Y) -- node[right] {\(f\)} (field);
  \draw[map] (X) -- node[below left=1pt] {\(F^*f\)} (field);
  \node (pullback) at (0,-13) {\(F^*f\)};
  \node (function) at (30,-13) {\(f\)};
  \draw[map] (function) -- node[above] {\(F^*\)} (pullback);
  \path (current bounding box.south) ++(0,-2mm);
\end{tikzpicture}

  \caption[Pullback by composition]{In the triangle, start in \(X\),
    apply \(F\), then evaluate \(f\). Below, \(F^*\) takes the function
    \(f\) on \(Y\) back to the function \(F^*f\) on \(X\).}
  \label{fig:linear-algebra-i:pullback-composition}
\end{marginfigure}

For a vector space \(V\), its linear dual \(V^*\) consists of the linear
functions \(V\to\mathbb F\). More generally, write \(\Lambda^k V^*\) for
the vector space of antisymmetric \(k\)-linear functions on \(V\).
Here \(k\) is the number of inputs; it need not equal \(\dim V\).

\begin{definition}
  Let \(L:V\to W\) be linear and let \(A\in\Lambda^k W^*\).
  The pullback of \(A\) by \(L\) is the function
  \[
    (L^*A)(v_1,\ldots,v_k)
      =A(L(v_1),\ldots,L(v_k)).
  \]
\end{definition}

Linearity of \(L\) preserves multilinearity in each input, and interchanging
two inputs interchanges the corresponding inputs of \(A\). Thus \(L^*A\)
is again antisymmetric and multilinear. The induced map
\[
  L^*:\Lambda^k W^*\longrightarrow\Lambda^k V^*
\]
is linear in \(A\).

\subsection{The space of top-degree forms}

Suppose \(\dim V=n\). The functions in \(\Lambda^n V^*\) are the
top-degree forms on \(V\); over \(\R\), a nonzero such function gives
an oriented volume function.

\begin{proposition}
  \label{thm:linear-algebra-i:top-degree-forms}
  The space \(\Lambda^n V^*\) has dimension one. If
  \(A\in\Lambda^n V^*\setminus\{0\}\), then every
  \(B\in\Lambda^n V^*\) is a unique scalar multiple of \(A\).
  For any ordered basis \((e_1,\ldots,e_n)\), this scalar is
  \[
    B=bA,\qquad b=\frac{B(e_1,\ldots,e_n)}{A(e_1,\ldots,e_n)}.
  \]
\end{proposition}

\begin{proof}
  The coordinate expansion in
  \autoref{eq:linear-algebra-i:antisymmetric-basis-expansion}
  determines each top-degree form by its value on the ordered basis.
  Since \(A\ne0\), its basis value cannot be zero. The displayed choice
  of \(b\) makes \(B\) and \(bA\) agree on that basis, so the same
  expansion gives \(B=bA\). Thus the single nonzero vector \(A\)
  is a basis of \(\Lambda^n V^*\).
\end{proof}
